% Aquí debes exponer las leyes de la física y los principios
% que rigen el comportamiento de la mezcla, independientemente de la sustancia
% computación que sustentan tu herramienta. 
% No defines términos, sino que explicas teorías.
% \section{El concepto de Estado en Termodinámica}


%%% OK

\subsection{Equilibrio Líquido-Vapor}

    Cuando las moléculas de un líquido tienen suficiente energía para escapar de la superficie, sucede un cambio de fase

\subsection{Presión de vapor}
    Cuando algunas moléculas de un líquido tienen energía suficiente para vender las fuerzas atractivas, y pasar al estado gaseoso, el proceso se llama evaporación o vaporización (Universidad de alcalá, cap 8)


\subsection{Mezclas Gas-Vapor}
Las mezclas debajo de la temperatura crítica


\subsection{Propiedades de los Fluidos}

  \figCoordinatePV{probando01}

  Esta el la figura \ref{probando01}

  \figSelectModelThermo{fig:selectEosModel}
  La selección de modelos \ref{fig:selectEosModel} \parencite{bib:DimianBildea2014}

% \subsubsection{Naturaleza de las mezclas Gas-Vapor}
% el
% \subsubsection{Naturaleza del Equilibrio}

% \subsubsection{Equilibrio Dinámico}
% \subsubsection{Sistema Termodinámico}

% \subsubsection*{Modelo para el Equilibrio Líquido-Vapor}
% \subsubsection*{Variables Macroscópicas Volumen, Presión, Temperatura}

% \subsubsection{Temperatura}
% En la actualidad se define en función a la constante de Boltzmann\footnote{Generado mediante \href{https://docs.scipy.org/doc/scipy/reference/constants.html}{cipy} y pythontex} {$k$}, \pcu{Boltzmann constant} en
% \subsubsection{Presión}
% En general, la presión se define como la fuerza normal ejercida sobre una unidad de area. Se habla de presión cuando el sistema consiste de fluidos (gas o líquido); pero cuando este es sólido, se usa el término \textit{esfuerzo normal} \textcite{YunusA.Cengel2015}

% \begin{equation}
%   \eqPressure
% \end{equation}

% \figMolecularColissions

% de un sistema en fase gaseosa, es la presión que aporta uno de los componentes


% \subsubsection{Fugacidad}
%     La fugacidad de un componente en una mezcla depende de la temperatura, presión y la composición de la mezcla \parencite[p. 248]{ReidPrausnitzPoling1987}

% \subsubsection{Presión Parcial}

% \subsubsection{Presión Parcial}
% como la fracción de uno de los componentes respecto a la suma de cada una de ellas, denominado presión total.

% \eqGasCompositionFromPressure
% Donde \symCompGV{i} es la composición molar, {\P{i}} es la presión del componente {${i}$}.

% la expresión \eqref{eqn:thermo:GasCompositionFromPressure}, también puede expresarse de la siguiente forma:

% \eqGasCompositionFromTotalPressure
% Donde % Implementar



% \subsubsection*{Ley de Henry}
%     En un sistema (L - V) en equilibrio a cierta temperatura, cuando un gas diluido se disuelve en un líquido y la fracción molar del gas tiende a cero; la presión en la fase gaseosa del componente diluido es directamente proporcional la composición en la fase líquida; se expresa de la siguiente manera:

%     \eqnHenrysLaw
%     \namingEqnHenrysLaw

%     De la ecuación ...

% \subsubsection*{Ley de Raoult}
% Se requiere hacer las siguientes suposiciones para las consideraciones para la ley de Raoult \parencite{Smith.VanNess.Abbott2007}
% \begin{itemize}
%   \item La fase de vapor es un \textbf{gas ideal}; significa que la aplicación se aplica solo a presiones bajas y moderadas.
%   \item la fase líquida es una \textbf{solución ideal}; implica que la validez aproximada se da cuando las especies que constituyen el sistema, sean químicamente semejantes\footnote{Especies moleculares de similar tamaño y de la misma naturaleza química.}
% \end{itemize}

% En base a las suposiciones mencionadas se deduce la siguiente expresión matemática

% \begin{equation}
%   \eqnRaoultLaw
% \end{equation}
% \namingEqnRaoultLaw

% % obs
% De la ecuación \eqref{eqn:thermo:RaoultLaw}, la expresión del lado izquierdo se conoce como, presión parcial del componente ${i}$.




% % El gas ideal permite la comparación del comportamiento del gas real 


% \begin{align}
%   f_{i}^{\degree, \text{ fase vap}} & = \hat{\phi_{i}} \symCompGV{i} \pressure \\
%   f_{i}^{\degree, \text{ fase liq}} & = \hat{\phi_{i}} \symCompL{i} p_{i}
% \end{align}

% \begin{flalign}
%   f_{i}^{\degree, \text{ fase vap}} & = \hat{\phi_{i}} \symCompGV{i} \pressure \\
%   f_{i}^{\degree, \text{ fase liq}} & = \hat{\phi_{i}} \symCompL{i} p_{i}
% \end{flalign}

% \begin{table}[htpb]
%   % \centering
%   \caption{Ley de Raoult y su versión modificada para la fase vapor y la fase líquida.}
  
%   \vspace{3mm}
%   \rowcolors{1}{gray!20}{}
%   \begin{tabular}{ccc}
%     \hline
%     \rowcolor{color_aquamarine} Variable & En la Ley de Raoult                              & En la Ley de Raoult modificada                               \\
%     \hline
%     $\hat{\phi_{i}}$                     & 1.00                                             & 1.00                                                         \\
%     \symCompGV{i}                        & \symCompGV{i}                                    & \symCompGV{i}                                                \\
%     $f_{i}^{\degree \text{, fase vap}}$  & \pressure                                             & \pressure                                                         \\
%     ${\gamma}_{i}$                       & 1.00                                             & ${{\gamma}_{i}}$                                             \\
%     $f_{i}^{\degree, \text{ fase liq}}$  & $p_{i}$                                          & $p_{i}$                                                      \\
%     Ec. resultante                       & $\symCompGV{i} = \frac{\symCompL{i} {P_{i}}}{P}$ & $\symCompGV{i} = \frac{ \gamma_{i} \symCompL{i} {P_{i}}}{P}$ \\ \hline
%   \end{tabular}
%   % Tomado de: PHYSICAL AND CHEMICAL EQUILIBRIUM FOR CHEMICAL ENGINEERS
  
%   \bigskip
%   \small\textit{Nota}. Tomado de  Physical and Chemical Equilibrium for Chemical Engineers, pág 345.
% \end{table}



% %-------------------------------------------------------------------
% Según \parencite{Engel2019}, una limitación que presenta la ley de los gases ideales es que no predice bajo las condiciones apropiadas la licuefacción\footnote{cambio de estado de una sustancia cuando pasa del estado gaseoso al líquido} de gases.


% \subsection{CAPACIDAD CALORÍFICA DE LA MEZCLA}
%     La capacidad calorífica de una mezcla de vapor-gas, se expresa en términos de la masa (mas no por mol) del gas; a menudo se supone constantes dentro de un intervalo de condiciones para los cálculos y esta manera la capacidad calorífica es:

%     \begin{equation}
%       \eqnHeatCapacityMixture
%     \end{equation}

%     % \namingImplHeatCapacity


% \section{ECUACIONES DE ESTADO}
%     Los estudios realizados por R. Boyle(1662), Charles () y Gay Lussac, sentaron los fundamentos para el estudio del comportamiento los gases. La generalización matemática del comportamiento de los gases es la siguiente:

%     \begin{equation}
%       \eqGeneralGases
%     \end{equation}
    
%     Donde {\zFactor} es el factor de compresibilidad, {\pressure} es la presión absoluta, {\Vmolar} es el volumen molar, {\temperatureAbsolute} es la temperatura y {\univConstGasEos} es la constante universal de los gases.


% \subsection{ECUACIÓN DE GAS IDEAL}
%     Un caso especial de la ecuación de estado los gases, es cuando el factor de compresibilidad (\zFactor) es igual a $1$. A esta forma de la ecuación general se llama \textit{ecuación ideal de los gases}, se expresa mediante la siguiente ecuación de estado:

%     \begin{equation}
%       \eqIdealGas
%     \end{equation}





%     Diversos autores señalan que esta ecuación (aunque existen otras más desarrolladas), es suficiente para describir el comportamiento de los gases, y que es una buena aproximación para gases reales a bajas presiones y altas temperaturas \parencite{YunusA.Cengel2015,ReidPrausnitzPoling1987,Smith.VanNess.Abbott2007}.
  
% \subsection{ECUACIÓN VIRIAL GENERAL}
%     \label{cap2:subsec:tagEqnvirial}

%     Según \parencite{Dymond2002}, muchas ecuaciones de estado han sido propuestas para representar el comportamiento $\pressure-\volume-\temperature$ de los gases reales desde el punto de vista teórico; de los cuales, la ecuación virial es el más satisfactorio.

%     Para \parencite[p. 37]{Sengers1987equations}, la ecuación virial de estado de los gases, es una de las ecuaciones que intenta describir de manera exacta y precisa el comportamiento de los fluidos reales. 

%     Una ventaja del uso de la ecuación virial es la existencia amplia de información experimental \Parencite[p. 38]{Sengers1987equations}, 

%     La ecuación virial puede ser apreciada como una serie de Maclaurin\footnotetext{También llamado serie de Taylor}, de ${\pressure}/{\univConstGasEos \temperatureAbsolute}$; en vista de su naturaleza polinomial, puede escribirse como serie de potencias de $1/\Vmolar$, el cual se muestra a continuación:

%     \begin{equation}
%       \eqVirialForVDevelop{}
%     \end{equation}
%     \namingEqVirialForVDevelop
    
%     El 2do y el 3er coeficiente virial ha sido estudiado por más de 100 años, todos ellos muestran que no existe correlación \parencite{thirdVirialTing}


%     % Se considera las siguientes suposiciones:

% \subsubsection*{Regla de Fase}  
%     La regla de fases fue desarrollada por 
%     \begin{equation}
%     \label{eqn:thermo:phase_rule}
%       \eqPhaseRule
%     \end{equation}

%     Donde,\newline 
%     \freeDegree : es el número de grados de libertad o las variables que se tienen que especificar, \newline 
%     \componentNumber : representa al número de componentes y \newline
%     \phaseNumber : es el número de fases (sólido, líquido o gaseoso) en el sistema.
%     % \includesvg[width=4.0in]{Imagenes/Vectorial/thermo/coord_P_V.svg}

% \subsubsection*{Dependencia de los coeficientes viriales de la temperatura}
%     \begin{equation}
%       \bCoeffDependentOfT{}
%     \end{equation}
    
%     \begin{equation}
%       \cCoeffDependentOfT{}
%     \end{equation}

% \subsubsection*{Dependencia de los coeficientes viriales de la composición}
%     Para un gas puro, de acuerdo a la ecuación \eqref{eqn:thermo:phase_rule} se necesita, dos grados de libertad (\freeDegree),  

%     Donde, $x = 1$ 

% \subsubsection{Métodos para la estimación de los coeficientes viriales}
%     Diversos autores señalan que los coeficientes viriales pueden derivarse de \parencite[p. 4.13]{Poling2001} de la teoría molecular
%     Para \parencite{Trusler10.1039/9781782627043-00152}, los métodos para la medición de los coeficientes viriales pueden clasificarse en:

%     \begin{enumerate}[label=(\alph*)]
%       \item Métodos directos
%         \begin{itemize}
%           \item Aparato de Burnet
%           \item Densímetro de doble plomada acoplados magnéticamente
%         \end{itemize}
%       \item  Métodos indirectos
%         \begin{itemize}
%           \item Flujo calorimétrico
%           \item Velocidad del sonido
%         \end{itemize}
%     \end{enumerate}

%     para \parencite{Dymond2002}, los métodos para la determinación de los coeficientes viriales pueden clasificarse de la siguiente manera:

%     \begin{itemize}
%       \item Medidas de {\pressure-\Vmolar-\temperatureAbsolute},
%       \item Medidas de la velocidad del sonido
%       \item Medidas de Joule-Thomson
%       \item Medidas de índice de refractividad y permitividad relativa
%       \item Medición de la entalpía de vaporización y la presión de vapor. 
%     \end{itemize}

% \subsubsection{Estimación del Segundo Coeficiente Virial}
  
%     El segundo coeficiente virial representa a la interacción entre pares de moléculas

% %%%%
% Son diversos los trabajos que se han realizado para determinar el segundo coeficiente virial. \autocite[postnote]{woolley1969second} trabajó para la formulación analítica del segundo coeficiente virial para una función esféricamente simétrica al cual denominó, \textit{potencial de par realista}. \parencite{pitzer1990second} también trabajó para estimar el segundo coeficientes virial para compuestos no polares y de baja polaridad. Para compuestos polares se han realizado diversos estudios para estimar el segundo coeficiente virial, entre los cuales destacan los estudios de \parencite{maris1985interaction}, \parencite{tarakad1977improved}, y entre otros estudios se han realizado . \parencite{vetere2007simple} propuso una modificacion a método de Pitzer para predecir el segundo coeficiente virial para compuestos puros.

% \parencite{klotz1985improved} presenta una mejora para el cálculo del 2do coeficiente virial



% Para la temperatura $\temperatureAbsolute$ y composición molar $\symCompGV{}$, ambos constantes; de la ecuación \eqref{eqn:virialB}, cuando $1 / \Vmolar\ \to 0$,  el segundo coeficiente virial $\bCoeffVirial{}$, puede escribirse como:

% \begin{equation}
%   \bCoeffVirial{} = \lim_{\frac{1}{\Vmolar}\ \to {0}} \Vmolar \left(\frac{\pressure \Vmolar} {\univConstGasEos \temperatureAbsolute} - 1 \right)
% \end{equation}

% Para \parencite[p. 473]{beattie1942second}, la expresión $\Vmolar \left({\pressure \Vmolar} / {\univConstGasEos \temperatureAbsolute} - 1 \right)$, puede ser calculado a partir de las mediciones de la compresibilidad isotérmica de la sustancia pura para su fase gaseosa; en caso de una mezcla, debe ser de composición molar constante; la compresibilidad isotérmica debe establecerse frente a $1 / \Vmolar$. El valor de la constante queda establecido en la intersección de la función con el eje $1/ \Vmolar$.

% Para una mezcla binaria, $\bCoeffVirial{}$ está relacionado con los coeficientes viriales de los componentes puros y sus composiciones, de la siguiente manera:

% \begin{equation}
%   \eqSecondVirialCoeffSum{}
% \end{equation}

% Para un sistema de dos componentes ($\symNumOfCmpnts = 2$),
% \begin{equation}
%   \eqSecondVirialCoeffDevelop{}
% \end{equation}
% \namingEqSecondVirialCoeff{}

% Conociendo los coeficientes viriales para los componentes puros, hace falta establecer las relaciones matemáticas para {\bCoeffVirial{\nameG \nameV}}.

% \subsubsection{Estimación del Tercer Coeficiente Virial}
% Para \parencite{Dymond2002}, los valores son preferentemente determinados mediante la información de compresibilidad de gases 

% \begin{equation}
%   \cCoeffVirial{} = \lim_{\frac{1}{\Vmolar}\ \to {0}} \left[ \Vmolar \left(\frac{\pressure \Vmolar} {\univConstGasEos \temperatureAbsolute} - 1 \right) - \bCoeffVirial{} \right] \Vmolar
% \end{equation}


% \begin{equation}
%   \eqThirdVirialCoeffSum{}
% \end{equation}

% Desarrollando para un sistema de dos componentes:

% \begin{equation}
%   \thirdVirialCoeffDevelop{}
% \end{equation}

% \namingEqThirdVirialCoeff{}

% \figVirialBCBehavior

% \subsection{Convergencia de la serie virial}

% \subsection{Teoría del Sonido}
% \subsubsection*{Velocidad del Sonido en los Gases}

% \parencite[p. 186]{Wilhelm2010}, define la velocidad del sonido como:

% \begin{align}
%   \label{cap:subsec:velocidadSonido}
%   \SoundSpeed{}
% \end{align}

% Donde $\pressure$  es la presión, $\density$ es la densidad, $\entropy$ es la entropía

% \subsection{Correlación, predicción y estimación de los coeficientes viriales}
% \label{cap2:subsec:tagEqnPredict}




% Según \parencite[ p. 12]{Dymond2002}, el método basado en la teoría del sonido, es el método que mejor se acerca a los valores reales; de esta manera, el segundo coeficiente virial queda expresado de la siguiente forma:
% \begin{equation}
%   \eqSndVirialCoeff{1}{2}{\temperatureAbsolute}
% \end{equation}


% Donde $f_{12}$ = $ \ce{exp}{\{-U(R, \omega_1 , \omega_2)/kt\}} - 1 $. Para una molécula lineal:

% Desarrollo de las ecuaciones para el cálculo de los coeficientes viriales fluidos polares y no polares





% \subsection{VOLUMEN MOLAR PARA UNA MEZCLA}
%     Para una mezcla gaseosa el volumen molar
%     \begin{equation}
%       \eqVirialForVDevelopForMixture
%     \end{equation}
%     \namingEqVirialForVDevelopMixture


% \begin{align}
%   \eqVirialFactorCorrection{}
% \end{align}
% \namingEqVirialFactorCorrection{}


% \subsection{ENERGÍA INTERNA}
% \parencite{Treybal.Mass.Transfer.1981} la energía interna $\internalEnergy$ de una sustancia es la energía total que reside en la sustancia debido a su movimiento y a la posición relativa de los átomos y moléculas que constituyen la sustancia. Los valores absolutos de la energía interna no son conocidas; sin embargo, su valor puede ser calculado respecto a un estado standard definido pra la sustancia.

% \subsection{ENTALPÍA PARA LA MEZCLA}
%   \begin{equation}
%     \eqEnthalpyForMixture
%   \end{equation}
%     \namingEqEnthalpyForMixture


% \subsection{ENTROPÍA PARA LA MEZCLA}
%   \begin{equation}
%     \eqEntropyForMixtureDiscrete
%   \end{equation}
% %%%%%%%%%%%%%%%%%%%%%%%%%%%%%%%%%%%%%%%%%%%%%%%%%%%%%%%%%%%%%%%%%%%%%%%%%%%%%%%%%%%%%%%%%%%%%%%%%%%%%%%%%%%%%%%%%%%%%%%%%%%%%%%%%%%%%%%%%%%%%%%



% \figSelectModelThermo{fig:SelectModelOfThermo}















% \section{Marco Teórico}
% \label{sec:marco_teorico}

%     \subsection{Fundamentos de las Ecuaciones de Estado}
%         Para \textcite{bib:thermo:goodstein2002states}, una ecuación de estado es toda relación entre la presión (\pressure[]{}) y temperatura (\temperature[]{}) de un sistema termodinámico en estado gaseoso.

        
%     Justificación del uso de la correlación de Wagner para sustancias que no sean agua, representar el comportamiento cerca del punto crítico.



%     \subsection{Termodinámica de Mezclas de Gases Ideales}
%         El comportamiento de la mezcla de gases se basa el la ley de Dalton

%     \subsection{Naturaleza del Equilibrio}
%         Según \textcite{bib:thermo:VanNess}, el equilibrio es una condición estática, donde no ocurre cambio alguna en las propiedades macroscópicas de un sistema en el tiempo.
%         Para \textcite{bib:thermo:goodstein2002states} un sistema en equilibrio, es cuando el estado inicial se vuelve irrelevante; lo que significa que, independientemente del mecanismo de cambio de estado, de la manera en que las partículas puedan redistribuir energía y cantidad de movimiento entre sí, todo recuerdo de cómo comenzó el sistema debe eventualmente quedar borrado por las infinitas distintas formas que las partículas puedan redistribuirse. Cuando un sistema alcanza esta condición, se dice que está en \textit{equilibrio}.

%         y \textcite{bib:thermo:gibbs1874} dice: "Para el equilibrio de cualquier sistema aislado es necesario y suficiente que en todas las posibles variaciones de el estado del sistema que no alteran su energía interna la variación de su entropía sea menor o igual que cero"; es decir, un sistema está en equilibrio cuando ya no puede reorganizarse de ninguna manera que lo haga más desordenado.

%     \subsection{Tipos del Equilibrio}
%         Estos son algunas formas de equilibrio:
%         \subsubsection{Equilibrio Térmico}
%             Cuando la temperatura es uniforme en todo el sistema.
%         \subsubsection{Equilibrio Mecánico}
%             Cuando la presión es uniforme y no hay fuerzas que desequilibren el sistema
%         \subsubsection{Equilibrio Químico}
%             La composición química no cambia;
%         \subsubsection{Equilibrio termodinámico}
%             Cuando el sistema alcanza las tres formas de equilibrio (térmico, mecánico, químico).

%         \subsection{Equilibrio de Fases}
%             Una fase es una región de un sistema con propiedades químicas y físicas uniformes; dentro de un sistema pueden coexistir diversas fases. Las transiciones de una fase a otra están reguladas por los principios termodinámicos.

%             \eqnClausiusClapeyron
%     \subsection{Termodinámica del Equilibrio de Fases}

%         % \genericPhaseDiagram



% \subsection{Fundamentos de la Termodinámica de Mezclas}
% \begin{itemize}
%     \item \textbf{Mezclas Gas-Vapor:} Comportamiento del aire seco y vapor de agua.
%     \item \textbf{Ley de Dalton:} Presiones parciales en mezclas binarias.
%     \item \textbf{Gases Ideales vs. Reales:} Justificación del modelo utilizado para el rango de operación.
% \end{itemize}

% \subsection{Teoría de la Psicrometría}
%     \begin{itemize}
%         \item \textbf{Propiedades Termodinámicas:} Definición matemática de Humedad Específica, Relativa y Entalpía.
%         \item \textbf{Ecuaciones de Estado para el Vapor de Agua:} Modelos de presión de saturación (Ecuación de Antoine o Goff-Gratch).
%         \item \textbf{Procesos Psicrométricos:} Calentamiento sensible, enfriamiento con deshumidificación y mezcla adiabática.
%     \end{itemize}

%     \subsection{Propiedades psicrométricas}
%         \begin{itemize}
%             \item Temperatura de bulbo seco
%             \item Temperatura de bulbo húmedo
%             \item Temperatura del punto de rocío
%             \item Humedad relativa
%             \item Relación de humedad
%             \item Entalpía específica
%             \item Volumen específico
%         \end{itemize}




%     \subsection{Entalpía específica del vapor}
        

% \subsection{Geometría de la Carta Psicrométrica}
% \begin{itemize}
%     \item \textbf{Sistemas de Coordenadas:} Transformación de valores termodinámicos a coordenadas cartesianas ($x, y$).
%     \item \textbf{Líneas de Saturación:} Deducción de la curva límite de la mezcla.
% \end{itemize}



